\documentclass[12pt,a4paper]{report}
\usepackage[utf8]{inputenc}
\usepackage{amsmath}
\usepackage{amsfonts}
\usepackage{amssymb}
\usepackage{amsthm}
\usepackage{hyperref}

\usepackage{multicol}
\usepackage{fancyhdr}
\usepackage[inline]{enumitem}
\usepackage{tikz}
\usepackage{tikz-cd}
\usetikzlibrary{calc}
\usetikzlibrary{shapes.geometric}
\usetikzlibrary{positioning}
\usepackage[margin=0.5in]{geometry}
\usepackage{xcolor}

\hypersetup{
    colorlinks=true,
    linkcolor=blue,
    filecolor=magenta,      
    urlcolor=cyan,
    pdftitle={Tensors},
    pdfpagemode=FullScreen,
    }

%\urlstyle{same}

\newcommand{\CLASSNAME}{Math 5050 -- Special Topics: Tensor Analysis}
\newcommand{\STUDENTNAME}{Paul Carmody}
\newcommand{\ASSIGNMENT}{Chapter 7: Elements of Linear Algebra in Tensor Notation }
\newcommand{\DUEDATE}{January 2026}
\newcommand{\PROFESSOR}{Professor Berchenko-Kogan}
\newcommand{\SEMESTER}{Spring 2026}
\newcommand{\SCHEDULE}{TBD}
\newcommand{\ROOM}{Remote}

\newcommand{\MMN}{M_{m\times n}}
\newcommand{\FF}{\mathcal{F}}

\pagestyle{fancy}
\fancyhf{}
\chead{ \fancyplain{}{\CLASSNAME} }
%\chead{ \fancyplain{}{\STUDENTNAME} }
\rhead{\thepage}

\fancyfoot{} % clear all footer fields
\fancyfoot[LE,RO]{\thepage}
\fancyfoot[LO,CE]{-------\\\ASSIGNMENT}
%\fancyfoot[CO,RE]{To: Dean A. Smith}
\input{../MyMacros}

\newcommand{\RED}[1]{\textcolor{red}{#1}}
\newcommand{\BLUE}[1]{\textcolor{blue}{#1}}

\newcommand{\SUPP}{\operatorname{supp}}
\newcommand{\CLOSURE}{\operatorname{closure of}}
\newcommand{\BD}{\operatorname{bd}}
\newcommand{\HHN}{\mathcal{H}^n}
\newcommand{\COFF}[3]{\Gamma^{#1}_{{#2}{#3}}}
\newcommand{\VK}[1]{\textbf{#1}}
\newcommand{\VKR}{\VK{R}}
\newcommand{\VKZ}{\VK{Z}}

\begin{document}

\begin{center}
	\Large{\CLASSNAME -- \SEMESTER} \\
	\large{ w/\PROFESSOR}
\end{center}
\begin{center}
	\STUDENTNAME \\
	\ASSIGNMENT -- \DUEDATE\\
\end{center} 

\textbf{Notes  --  Linear Algebra vs Tensors}

\begin{tabular}{|c|c|c|}
	\hline
	& Linear Algebra & Tensors \\
	\hline
	Vectors & $\VK{v} = v^i\VK{e}_i$ & $\VK{V} = V^iZ_i$ \\
	\hline
	Matrix & $A=[ A_{ij} ]$ & $A_{ij}, A^i_j, A^{ij}$ \\
	\hline
	Basis & $\VK{E}=\COLVECTOR{\VK{e}_1\\\dots\\\VK{e}_n}$ : $\VK{E}'=\COLVECTOR{\VK{e}_{1'}\\\dots\\\VK{e}_{n'}}$ & $\VK{Z}_i = \PART{\VK{R}}{Z^i}$ : $\VK{Z}i' = \PART{\VK{R}}{Z^{i'}}$\\
	\hline
	& $\VK{v} = v^T\VK{E}$ & $\VK{V} = V^i\VK{Z}_i$ \\
	\hline
	Jacobian & $X$ & $J^i_j$ \\
	\hline
	Transformation & $ v' = X^{-1}v$ : $v=Xv'$ & $v^{i'} = J^{i'}_iv^i$ : $v^{i'} = J^{i}_{i'}v^{i'}$ \\
	\hline
	Inner Product & $(\VK{u},\VK{v}) = (\VK{e}_i,\VK{e}_j)u^iv^j$ & $\VK{V}\cdot\VK{U} = V^iU^i$ \\
	\hline
	Gram Matrix/Metric Tensor & $M = [ \VK{e}_i\cdot \VK{e}_j ]$ & $Z_{ij} = Z^iZ^j$ \\
	\hline
	& $(\VK{u},\VK{v}) = M_{ij}u^iv^j$ & $\VK{Z}_i\cdot\VK{Z}_j = Z_{ij}Z^iZ^j$\\
	\hline 
	& $(\VK{u},\VK{v}) = u^TMv$ & \\
	\hline
	& $M' = X^TMX$ & $Z_{i'j'} = J^i_{i'}Z_{ij}J^j_{j'}$\\
	\hline
	& $A' = X^{-1}AX$ & $A^{i'}_{j'} = A^{i}_{j} J^{i'}_iJ^j_{j'}$\\
	\hline
\end{tabular}\\


\HLINE
\textbf{Exercises}
\begin{enumerate}[label=\textbf{\arabic*.}]
\setcounter{enumi}{106}
\item Express the symmetric property
\begin{align*}
	A^T = A
\end{align*} in the tensor notation.

\item Express the product $C=A^TB$ in the tensor notation.

\item Express the product $C^T=A^TB$ in the tensor notation.

\item Express the product $C^T=A^TB^T$ in the tensor notation and explain why it follows that $(AB)^T = B^TA^T$.  This is an example of a theorem in matrix algebra that is best proved in the tensor notation.

\item Use the matrix notation to prove that, for inveritble matrices $A, B,$ and $C$,
\begin{align*}
	(ABC)^{_1}=C^{-1}B^{-1}A^{-1}
\end{align*}

\BLUE{\begin{align*}
`	\LET D &= BC \\
	(ABC)^{_1} &= (AD)^{-1} = D^{-1}A = (BC)^{-1}A^{-1} = C^{-1}B^{-1}A^{-1}
\end{align*}
}

\item Linear algebra offers a derivation that does not require claculus.  Introduce a new vector $r = A^{-1}x-b$ and express $f(x)$ in terms of $r$.  Show that the minimum of the new quadratic form occurs at $r=0$.

\item Without the tensor calculus, the calculus approach proves surprisingly cumbersome.  Derive the partial derivatives of 
\begin{align*}
	f(x,y) &= \frac{1}{2}A_{11}x^2+!_{12}xy+\frac{1}{2}A_{22}y^2 -b_1x-b_2y
\end{align*}

\item Perform the same task in three dimensions for $f(x,y,z)$.

\item Show that the Hessian $\partial^2 f/\partial x^i\partial x^j$ equals 
\begin{align*}
	\frac{\partial^2 f}{\partial x^i\partial x^j} = A_{ij}
\end{align*}

\item Show that equation (7.71) can be written as 
\begin{align*}
	A_{ij}x^j=\lambda x_i
\end{align*}and as
\begin{align*}
	A^i_j x^j = \lambda x^j
\end{align*}

\item Show that the eigenvalues of the generalized equation (7.71) are invariant with respect to change of basis.

\item Show that the eigenvalues of the generalized equation (7.71) are given by the Rayliegh quotient
\begin{align*}
	\lambda = A_{ij}x^ix^j
\end{align*}

\end{enumerate}

\end{document}